\scanpage{118}

% Two long oblique pen strokes traverse the typed draft.
\deleted{Page 33. Ligne $-9$ à $-1$, remplacer le texte par le suivant:}

\deleted{\handwritten{«} \uline{Remarque}s7.10. a) Lorsque $G$ est à fibres connexes, alors on peut
dans c) abandonner l'hypothèse que les $\deleted{\unsure{p}}\,{}_{p^{\handwritten{\unsure{\nu|v}}}}H'$ soient séparés. En
effet, avec les réductions faites dans la démonstration, on peut
supposer $u$ bijectif, donc $H$ à fibres connexes, donc $H$ séparé en vertu du théorème de séparation de Raynaud (Exp $\mathrm{VI}_{\mathrm B}$).}

\deleted{b) D'autre part, des exemples de Raynaud montrent qu'on
ne peut en général dans c) abandonner ni l'hypothèse que les $ {}_{p^{\handwritten{\nu}}}G'$
soient finis sur $S'$, ni que les $ {}_{p^{\handwritten{\nu}}}H'$ soient séparés sur $S'$.}

\deleted{c) \inserted{On peut montrer que} \deleted{L}la conclusion de a) est encore valable si, au lieu de supposer $G$ à fibres connexes, on le suppose de type multiplicatif; \deleted{\unsure{i.e.}}
tout monomorphisme $u:G\handwritten{\longrightarrow}H$ de $S$-préschémas en groupes\deleted{\unsure{,}} de présentation finie, avec $G$ de type multiplicatif, \deleted{\unsure{et H|et N}} est une immersion. \handwritten{»}}

\deleted{Ligne 14, lire $S'$ au lieu de $S$.}

\handwritten{Page 35, ligne 9,10, supprimer ``lisse sur $k$''.}

% archive sheet number: 117
